\section{Discussion, Limitations, and Reproducibility}
\label{sec:discussion}
\paragraph{What the structural profiles establish.}
The common-size contrast between outward stars and incoming trees shows that task-family difficulty is not explained by node count alone. Size curves then show that increasing reference size need not produce the same profile across families or evaluated models. A useful evaluation should therefore stratify by family, model, and evidence condition rather than collapse them into a universal graph-difficulty score. This conclusion concerns the measured entity-selection task; it does not rank the general reasoning ability of a model.

The structural axis remains conditional. Incoming/outgoing orientation, traversal depth, terminal-type rules, entity/type distributions, and question wording vary with family. Accepted samples can differ across sizes. Fixed candidate count does not fix target prevalence, context length, graph density, or per-entity selection difficulty. Ceiling effects obscure small losses in strong conditions. The study does not isolate topology from all these factors, and we report no regression-adjusted causal or significance claim for motif or size.

\paragraph{Relation to the KG appendix.}
The KG experiment in Tool Graph Complexity already studies motif and size effects with explicit relations and adjacent neighbours~\cite{anonymousToolGraph}. We position our work as a focused examination of that related KG line under additional evidence and presentation contrasts. We do not claim independent validation of Appendix~O or new construction of its KG substrate. Its exact historical row overlap with this archive is not independently quantified. Scores across the papers should not be directly subtracted or interpreted as changes in a model's capability. The additional value here is the explicit conditional analysis and interpretation, not rediscovery of KG motif or distractor effects.

\paragraph{Reference validity and downstream utility.}
Reference coverage, minimal graph support, indispensable visible cards, and correct answers are distinct. Matching $W$ demonstrates adherence to the route-coverage convention. A smaller selected set could still contain sufficient evidence for an endpoint answer; a large set can receive nonzero F1 without grounded reasoning. Analytical references illustrate this: select-all F1 is 0.364 on Core and 0.405 on Full Raw, despite not using the question or facts. We provide these as metric reference points, not executed retrieval systems or competitive QA baselines. There is no validated downstream QA improvement in this study.

\paragraph{Sampling, uncertainty, and generalization.}
One corporate KG and its accepted frozen tasks cannot establish transfer to other domains, snapshots, model versions, or prompts. Anchors repeat and reference entities cross anchor clusters. Shared entities connect 578 of 700 Core tasks and 3,265 of 3,400 Full Raw tasks in their largest reference-overlap components. Anchor bootstrap intervals thus rely on cross-anchor independence conditional on the fixed KG; they do not control arbitrary shared-entity dependence. The nominal simultaneous bands and task-bootstrap sensitivity should be read within that restriction. Transport retries and multiple response stages are not independent repetitions, and we have no generation-seed replication study.

\paragraph{Historical evidence and reproducibility.}
The archived observation, input, response, and task-specification files retain IDs and hashes, allowing task/reference checks and bounded offline score aggregation. Some Gemma outcomes have cached scores without raw replies; unavailable call settings cannot be inferred from current defaults. Incomplete incremental histories prevent equal-budget or mechanism claims about interactive access. The full source snapshot was not loaded during the audits, so frozen-context agreement does not prove completeness against the original database.

The accompanying analysis materials identify source rows, planned denominators, pairing keys, failure policy, resampling units, and figure provenance. Motif/size plots are derived from existing audited summaries and add no new inference. Controlled and legacy parsing/scoring conventions remain distinct; legacy outcomes are not pooled into the main results. Appendix queries document read-only access/checks without representing their execution as an experiment. Public release of raw resources and the resolution of any access restrictions are separate from the reproducibility of these archived summaries.

\paragraph{Data and writing responsibility.}
The analysis uses existing corporate-KG experimental artifacts and includes no new human-subject study or new model calls. We do not infer consent, redistribution rights, or an ethics-board decision from the presence of an archive. Generative AI assistance was used to organize the evidence audit, prepare analysis/plotting code, and draft and edit this manuscript; the numerical summaries derive from the saved experimental data rather than invented model outputs. The authors remain responsible for checking the manuscript and its relationship to the related KG study.
